\begin{abstract}
Condition monitoring sets thresholds on scalar indicators, and those
thresholds transfer poorly between machines because the same health state
produces different values under different mounting, load and sensor placement.
This paper formulates diagnosis instead as inverse recovery over a quotient
space induced by the nuisance transformations of the observation process, in
which the object to be recovered is an equivalence class rather than a point.
Three consequences are established. Under bounded nuisance the nominal orbit
is relatively compact, so its image under a continuous representation has a
finite radius and there is a well-defined quantity to estimate rather than a
threshold to tune. A correct diagnostic decision factors through the quotient,
which licenses one shared representation in place of a per-asset classifier.
And that radius is exactly a distribution-free tolerance limit on the nominal
residual, so its attained coverage is known exactly as \(k_\alpha/(n+1)\)
rather than bounded: at a target of \(98\%\), \(49\) nominal observations
suffice to certify a boundary, and a degenerate range above that floor is
identified in which the certified estimator is the sample maximum and can
appear to have converged while remaining maximally variable. Validation
proceeds in three tiers. Exhaustive verification reproduces the exact coverage
to within \(\satMaxZ\) Monte Carlo standard errors across six deliberately
dissimilar residual distributions, and isolates the cost of the interpolated
estimator used in practice at \(1.9\) percentage points of coverage. On a
bearing simulator, one nominal-only operator detects three fault families at
\(\etoeDetectLo\)--\(\etoeDetectHi\%\) at a \(\etoeFAR\%\)
false-alarm rate, and a probe shows that coverage alone cannot certify a
representation: two operators calibrated identically, at the same attained
coverage, differ by more than two orders of magnitude in what they accept. On
CWRU, calibration from roughly fifty sequential observations matches a
reference fitted and calibrated on each speed's complete healthy record; on
\(\nasaBearings\) NASA IMS run-to-failure bearings, radii commissioned from
\(\nasaNLo\)--\(\nasaNHi\) observations and never adapted give persistent
departure with \(\nasaLeadLo\) to \(\nasaLeadHi\) hours of warning, with
no fault labels at any stage. Repeating the field pipeline across training
seeds shows the radius moving by \(4\)--\(7\%\) while the same radius
normalised by each asset's own residual level moves by under \(1.3\%\), so
the diagnostic content lies in the geometry of the nominal region relative to
its own scale rather than in the absolute value of the residual. A
reproducibility package regenerates every table, figure and quoted number.

\end{abstract}

\section{Introduction}
\label{sec:introduction}

A rotating machine in steady operation does not produce arbitrary signals. Its
vibration is the output of a mechanism with finitely many degrees of freedom,
observed through a sensing chain that introduces its own bounded distortions,
and the set of waveforms it can emit while healthy is correspondingly
restricted. Load, speed, mounting, sensor gain and acquisition phase move the
observed signal about, but they move it within a region rather than across an
unbounded space.

This paper takes that restriction as its governing premise and states it
sharply, because it is the premise that makes the rest of the construction
either meaningful or vacuous: \emph{physically generated signals do not
exhibit unbounded or adversarial variation, and departure from a compact
nominal region is what an anomaly is}. The claim is not that nominal signals
are often well behaved. It is that the nominal set is relatively compact, and
therefore has a finite extent that can be estimated from finitely many
observations. If that fails, no finite boundary exists, and a method built on
estimating one is not merely inaccurate but meaningless.

Stating the premise this way makes it refutable, and Section~\ref{sec:results}
is arranged so that it could have been refuted. A process whose nominal
variation did not concentrate would show a boundary estimate that kept growing
with the number of observations rather than settling; the experiments measure
exactly that quantity, and report both the regime in which it settles and the
regime in which it does not.

\paragraph{Diagnosis as inverse recovery.}
An observation is the output of a forward process. A vibration window \(x\) is
produced by a mechanism in some health state \(\theta\), acted on by nuisance
factors \(\eta\) that carry no diagnostic information, so that
\(x=F(\theta,\eta)\). Recovering \(\theta\) from \(x\) is an inverse problem,
and in the classical sense an ill-posed one: the forward map is not injective,
because distinct \((\theta,\eta)\) pairs can produce the same observation, and
its inverse is unstable where it exists.

The usual response is to identify the forward model and invert it. For
rotating machinery this is neither tractable nor necessary. It is intractable
because a faithful forward model of a bearing in an installed drivetrain
requires contact mechanics, lubrication state, structural transmission paths
and sensor dynamics that are not identifiable from routine monitoring data. It
is unnecessary because diagnosis does not require \(\theta\) itself. It
requires \(\theta\) only up to the nuisance factors: two windows differing by
mounting phase or sensor gain must receive the same diagnostic verdict, so
the object to be recovered is not a point but an equivalence class.

That observation is the central claim of this work, and we state it in the
form used in the companion technical note~\citep{disanti2026quotient}:
diagnostic monitoring can be formulated as inverse recovery over a quotient
space induced by nuisance transformations of the observation process.

The reformulation changes what must be estimated. In the quotient space
\(\quotient\), the nuisance orbit of a nominal signal collapses to a single
element, the nominal set becomes a region rather than a scattered cloud, and
the diagnostic question becomes membership in that region. Three things then
follow that do not follow in the ambient space. Membership is decidable from
nominal data alone, so no fault examples are needed. The region has a finite
extent, because bounded nuisance makes the orbit relatively compact. And that
extent is a scalar, which is exactly the kind of quantity for which
distribution-free statistical guarantees exist.

\paragraph{What the premise buys, quantitatively.}
The last point is where a qualitative geometric argument becomes an
engineering statement. If the nominal region has a finite radius, then
estimating that radius is a tolerance-limit problem, and the coverage of the
estimate can be certified without assuming a distribution for the residual.
The certification is not asymptotic and does not appeal to a limit: at a
target coverage of \(98\%\), forty-nine nominal observations suffice, and the
attained coverage is known exactly rather than bounded.

That number is the operational form of the premise. A regime that can be
bounded from tens of observations is compact in the sense the premise asserts;
a process without that structure would require the estimate to keep growing.
The same analysis also delimits its own validity, and this is worth stating in
the introduction rather than burying in the results: below forty-nine
observations no boundary is certifiable at all, and between forty-nine and
ninety-eight the certified estimator is the sample maximum, so it can appear
to have converged while remaining maximally variable. A commissioning length
chosen just above the floor satisfies the guarantee and is nonetheless a poor
choice.

\paragraph{Contributions.}
\begin{enumerate}
\item A quotient-space formalization of anomaly detection, with two results
      that carry the argument: the nominal orbit is relatively compact under
      bounded nuisance (Proposition~\ref{prop:compactness}), and a correct
      diagnostic decision factors through the quotient
      (Proposition~\ref{prop:quotient_diagnosis}), which is what licenses
      seeking one shared representation rather than a per-asset classifier.

\item A distribution-free coverage guarantee for the locally estimated
      operational radius (Proposition~\ref{prop:coverage}), given with proof.
      The statistical machinery is classical; the contribution is recognizing
      that the per-asset diagnostic radius is exactly a distribution-free
      tolerance limit on the nominal residual. Two consequences follow that
      are specific to the diagnostic setting and, to our knowledge, have not
      been stated for it: a minimum commissioning length
      (Corollary~\ref{cor:min_commissioning}), and a degenerate range above
      that floor in which the estimator is the sample maximum
      (Corollary~\ref{cor:degenerate_range}).

\item A saturation criterion connecting this formulation to the
      functional--topological account of~\citet{disanti2026blueprints}, with
      an explicit statement of what rapid convergence means quantitatively
      (Definition~\ref{def:saturation_radius},
      Proposition~\ref{prop:radius_consistency}), together with an operator
      adequacy criterion showing that coverage alone cannot certify a
      representation: two operators calibrated identically, at the same
      attained coverage, differ by more than two orders of magnitude in what
      they accept.

\item Validation across three tiers of increasing realism: synthetic assets
      under controlled nuisance, CWRU under four operating regimes with an
      oracle-versus-online calibration ablation, and NASA IMS run-to-failure
      bearings commissioned independently and with no fault labels at any
      stage.
\end{enumerate}

\paragraph{Scope.}
Two boundaries should be set at the outset. The framework addresses membership
in a nominal region, not fault classification: it reports that an asset has
left its region, not which failure mode it has entered, and the taxonomy in
Section~\ref{sec:emergent_classification} grows from expert feedback rather
than from the geometry alone. And the guarantee is conditional on
exchangeability of the nominal commissioning residuals
(Assumption~\ref{ass:exchangeability}), which is an assumption about the
commissioning window and not a property of the method; Section~\ref{sec:results}
reports a test of it rather than presuming it holds.

